\chapter{Introduction}

\section{Background}

% This is a citation \citep{van2019sustainable}. Citations should be made using the \texttt{natbib} package, using the \texttt{citep} and \texttt{citet} commands.

The volume of scientific literature published globally continues to expand at an unprecedented rate. According to recent bibliometric analyses, the number of peer-reviewed articles published annually has grown exponentially over the past two decades, with major indexing services such as PubMed, IEEE Xplore, and arXiv collectively hosting millions of new preprints and journal publications each year \citep{bornmann2015growth}. Researchers, postgraduate students, and academic professionals across virtually all technical disciplines now face an overwhelming corpus of papers, conference proceedings, and technical reports that cannot feasibly be read in their entirety during a single literature review cycle. This information overload introduces significant cognitive burden, productivity bottlenecks, and increased risk of overlooking relevant prior work.
\newline \newline Traditional approaches to literature review manual reading, keyword-based database searches, and citation chaining remain effective for targeted queries but scale poorly as publication volumes increase. Conventional keyword matching algorithms, such as TF-IDF and BM25, rely on literal token overlap between queries and documents, fundamentally limiting their ability to capture semantic synonyms, paraphrases, or conceptual relationships expressed in different vocabulary \citep{robertson2009probabilistic}. A researcher searching for "methodology employed in the study" would fail to retrieve passages describing "experimental setup" or "algorithmic implementation" unless those exact terms appeared in the text. This lexical gap between how researchers formulate questions and how authors express ideas in papers represents a fundamental limitation of traditional information retrieval.
\newline \newline The emergence of transformer-based deep learning architectures \citep{vaswani2017attention} and pre-trained \gls{llm}s has fundamentally transformed the natural language processing landscape. Models such as Meta's LLaMA \citep{touvron2023llama}, Alibaba Cloud's Qwen \citep{yang2024qwen2}, and OpenAI's GPT series demonstrate remarkable capabilities in text comprehension, summarisation, and question answering across diverse domains. However, deploying \gls{llm}s directly for academic question answering introduces a critical reliability concern: hallucination. When prompted with questions about papers they have not been trained on, \gls{llm}s frequently generate plausible-sounding but factually incorrect or entirely fabricated responses \citep{ji2023survey}. In academic contexts, where factual precision and source traceability are paramount, hallucinated outputs can mislead researchers and compromise the integrity of literature reviews.
\newline \newline \gls{rag} has emerged as the leading architectural paradigm for addressing \gls{llm} hallucination in knowledge-intensive tasks \citep{lewis2020retrieval}. Rather than relying solely on the parametric knowledge stored within a language model's weights, \gls{rag} systems first retrieve relevant passages from an external knowledge base and then inject these passages as grounding context into the generation prompt. By constraining the \gls{llm} to answer based strictly on retrieved source material, \gls{rag} architectures substantially reduce hallucination while enabling source citations that allow users to verify and trace every claim back to the original document. This combination of semantic retrieval accuracy and generative fluency positions \gls{rag} as a particularly promising architecture for building research paper analysis tools.
\newline \newline Recent advances in dense vector representation learning, specifically Sentence-\gls{bert} \citep{reimers2019sentence} and lightweight variants such as all-MiniLM-L6-v2, further enhance retrieval quality by projecting text into continuous semantic vector spaces where geometric distance directly measures conceptual similarity. When combined with cloud-native vector databases such as Pinecone, which provide scalable high-dimensional indexing with sub-second query latency, these embedding models enable real-time semantic search across large document collections. Furthermore, contrastive fine-tuning techniques using objectives such as \gls{mnrl} \citep{henderson2017efficient} allow domain adaptation of general-purpose embedders to specialised corpora such as scientific papers significantly improving retrieval precision for academic question answering.
\newline \newline It is within this technological context that InsightPaper AI is conceived: a full-stack, end-to-end \gls{rag} system designed specifically for grounded research paper analysis, combining fine-tuned semantic retrieval, multi-model \gls{llm} generation, and an interactive evaluation framework.
\newline \newline Below as shown in Figure \ref{fig:fig01} is a graph illustrating the growth in annual scientific publications over the past 20-24 years, highlighting the increasing volume of research literature and the resulting information overload.


\begin{figure}[H]

\includegraphics[width=0.9\linewidth]{figures/1.1.png}

\caption{Growth of Scientific Publications}

\label{fig:fig01}

\end{figure}

\section{Problem Statement}

Despite the transformative potential of \gls{llm}s and semantic search technologies, several critical challenges persist in applying these tools effectively to academic research paper analysis:

\subsection{Hallucination and factual unreliability} General-purpose \gls{llm}s, when prompted with questions about specific research papers, frequently generate responses that contain fabricated claims, incorrect numerical results, or misattributed findings. Without access to the actual paper content during generation, the model relies on probabilistic patterns learned during pre-training, which may not accurately reflect the specific paper under analysis. This lack of grounding makes direct \gls{llm} querying unsuitable for rigorous academic work where precision and verifiability are essential.

\subsection{Lexical gap in retrieval}

Traditional keyword-based search systems fail to bridge the vocabulary mismatch between how researchers formulate questions and how academic authors express ideas. Dense semantic retrieval using pre-trained sentence transformers offers a solution, but general-purpose embedding models trained on broad web corpora may not adequately capture the specialised vocabulary, structural conventions, and discourse patterns characteristic of scientific literature. Domain-specific fine-tuning is required to close this gap, yet the effectiveness of such fine-tuning for \gls{rag}-based academic question answering remains underexplored.

\subsection{Lack of integrated evaluation frameworks}

Existing \gls{rag} prototypes and commercial tools for paper analysis rarely provide built-in mechanisms for quantitatively evaluating answer quality. Without automated metrics comparing generated answers to gold-standard references, users cannot objectively assess system reliability. A robust evaluation framework encompassing token-level accuracy, semantic alignment, and holistic quality judgement is needed.

\subsection{Limited multi-model flexibility}

Most existing systems are locked to a single \gls{llm}, preventing researchers from comparing response quality, inference latency, and reasoning depth across different model architectures and parameter scales. The ability to switch between models,for example, between a high-capacity reasoning model and a lightweight, low-latency model would provide valuable flexibility for different use cases.
\newline \newline These challenges collectively motivate the development of InsightPaper AI: a system that integrates fine-tuned dense retrieval, multi-model \gls{llm} generation with source grounding, and a comprehensive automated evaluation framework within a single, accessible web application.

\section{Aim and Objectives}


% \begin{figure}[H]

%   \includegraphics[width=0.9\linewidth]{figures/iris.png}

%   \caption{My figure}

%   \label{fig:fig01}

% \end{figure}

The overarching aim of this project is to design, implement, and evaluate InsightPaper AI, a \gls{rag} system that enables researchers to upload scientific papers as \gls{pdf}s and receive accurate, source-grounded answers to natural language questions, supported by verifiable source citations and automated quality evaluation. To achieve this aim, the following specific objectives are defined:

\begin{itemize}

\item To develop an end-to-end \gls{rag} system for uploading academic papers and generating source-grounded answers with citations.

\item To fine-tune the embedding model to improve retrieval performance on scientific papers and evaluate it using Recall@K and \gls{mrr}.

\item To integrate multiple large language models through the Groq \gls{api}, with user-selectable models and an automatic fallback mechanism.

\item To evaluate the system using token-level F1 score, semantic similarity, and \gls{llm}-as-a-Judge metrics for faithfulness, relevance, and completeness.

\item To build a user-friendly Streamlit web application that supports \gls{pdf} upload, paper summarisation, semantic search, figure extraction, and chat-based question answering.

\end{itemize}

\section{Research Questions}

This project is guided by one primary research question and four supporting sub-questions (5 research questions in total):

\begin{itemize}
    \item \textbf{Primary Research Question (RQ):} How effectively can an end-to-end \gls{rag} architecture, utilising a fine-tuned Sentence-Transformer embedder, a Pinecone vector index, and \gls{llm}s, retrieve relevant paper contexts and generate accurate, source-grounded answers for academic research papers?
    \item \textbf{Sub-Question 1 (SQ1 — Retrieval Adaptation):} To what extent does fine-tuning a lightweight embedding model (all-MiniLM-L6-v2) on SQuAD 2.0 and the Kaggle Summarized Research Papers dataset via \gls{mnrl} improve retrieval recall and mean reciprocal rank compared to the base pre-trained model?
    \item \textbf{Sub-Question 2 (SQ2 — Model Scaling \& Performance):} How do token-level F1 scores, dense semantic similarity metrics, and \gls{llm}-as-a-Judge ratings compare when evaluating Qwen 3.6 27B and LLaMA 3.1 8B model-generated answers against ground-truth dataset answers?
    \item \textbf{Sub-Question 3 (SQ3 — Multi-Metric Evaluation Framework):} How effective are token-level F1 score, dense semantic similarity, and multi-dimensional \gls{llm}-as-a-Judge ratings in evaluating answer accuracy, factual faithfulness, and explanation completeness?
    \item \textbf{Sub-Question 4 (SQ4 — Metadata Filtering \& System Latency):} How effectively can source-type metadata filtering (\gls{pdf} vs. Dataset) isolate document domain contexts while maintaining real-time query latency (under 1.5 seconds)?
\end{itemize}

\section{Thesis Structure}

This dissertation is structured into seven logical chapters, moving from theoretical foundations and architectural design to empirical implementation, fine-tuning, and rigorous multi-metric evaluation:

\begin{itemize}
    \item \textbf{Chapter 1 — Introduction:} Establishes the research background, problem statement, research aim and objectives, core research questions, and the structural overview of the dissertation.
    \item \textbf{Chapter 2 — Literature Review:} Evaluates foundational literature across deep learning document parsing, dense vector representations, contrastive fine-tuning, cloud-native vector databases, \gls{rag} architecture patterns, and \gls{llm} evaluation frameworks, identifying the empirical research gaps addressed by this work.
    \item \textbf{Chapter 3 — System Design and Architecture:} Details the hybrid Iterative Prototyping \gls{sdlc} methodology, functional and non-functional requirements, \gls{uml} Use Case specifications, high-level \gls{mvc} system architecture, sliding-window chunking parameters, Pinecone vector index schema, and multi-dimensional evaluation matrix design.
    \item \textbf{Chapter 4 — System Implementation and Prototypes:} Presents the concrete realization of InsightPaper AI across PyMuPDF parsing integration, Pinecone \gls{rag} indexing and retrieval pipelines, Groq \gls{api} \gls{llm} routing middleware (Qwen 3.6 27B and LLaMA 3.1 8B), four annotated \gls{ui} interface prototypes, Streamlit session state management, and the automated dataset evaluation scoring pipeline.
    \item \textbf{Chapter 5 — Embedding Model Fine-Tuning:} Documents the domain-adaptive fine-tuning lifecycle of \texttt{all-MiniLM-L6-v2} using \gls{mnrl} on 120 training papers from the Kaggle Summarized Research Papers dataset, detailing loss mechanics, training hyperparameters, and post-training Recall@5 and \gls{mrr} evaluation results.
    \item \textbf{Chapter 6 — Results and Comparative Evaluation:} Conducts quantitative empirical benchmarking across 30 held-out evaluation papers, analyzing retrieval performance gains, token-level F1 scores, 384-dimensional dense semantic similarity, multi-dimensional \gls{llm}-as-a-Judge ratings, and end-to-end system query latency.
    \item \textbf{Chapter 7 — Conclusion and Future Work:} Summarises the core research contributions, explicitly answers the primary research question and sub-questions, critically reviews system limitations, and outlines five strategic directions for future research.
\end{itemize}
